\section{Two-ended packets on the actual short segments}
\label{sec:v95-packets}
An inlet trace alone can overestimate a short pulse. Its other physical
end gives a second trace on the same incident component. We retain both
before summing the original source. In this section the endpoints are
physical endpoints $a_\sigma,b_\sigma$, not artificial prefix cuts.
Put $E_\varepsilon=1+\varepsilon^{-1}$ and
\[
 F_\sigma^-=f_\sigma^\varepsilon(a_\sigma+),\qquad
 F_\sigma^+=f_\sigma^\varepsilon(b_\sigma-),\qquad
 M_\sigma=\min\{F_\sigma^-,F_\sigma^+\}.
\]
For fixed $\rho,\delta>0$ let
\[
 \mathcal D_{\rho,\delta}(t)
 =\{\sigma:a_\sigma<t<b_\sigma,\ z_\sigma^-\ge\rho,
                                      \ 0<\ell_\sigma\le\delta\}.
\]
Integrating over this same family defines four nonnegative functions:
\begin{align}
 Q_{\rho,\delta}^\varepsilon(t)&=\int_{\mathcal D_{\rho,\delta}(t)}
                                      f_\sigma^\varepsilon(t)\,d\sigma,\notag\\
 I_{\rho,\delta}^\varepsilon(t)&=\int_{\mathcal D_{\rho,\delta}(t)}F_\sigma^-\,d\sigma,
 \quad O_{\rho,\delta}^\varepsilon(t)=\int_{\mathcal D_{\rho,\delta}(t)}F_\sigma^+\,d\sigma,
 \quad M_{\rho,\delta}^\varepsilon(t)=\int_{\mathcal D_{\rho,\delta}(t)}M_\sigma\,d\sigma.
 \label{eq:v95-packets}
\end{align}
In particular $Q$ is a restriction of the original weighted collision
source. The minimum in $M$ is taken on each physical segment, not
between two projected totals, and no sign cancellation is assumed.

\begin{theorem}[A complete two-ended replacement]
\label{thm:v95-packets}
At every finite count and almost every roof,
\begin{align}
 |Q_{\rho,\delta}^\varepsilon-M_{\rho,\delta}^\varepsilon|(t)
 &\le V_{\rho,\delta}^\varepsilon(t),\notag\\
 |I_{\rho,\delta}^\varepsilon-M_{\rho,\delta}^\varepsilon|(t)
 +|O_{\rho,\delta}^\varepsilon-M_{\rho,\delta}^\varepsilon|(t)
 &\le V_{\rho,\delta}^\varepsilon(t),\label{eq:v95-packet-error}\\
 V_{\rho,\delta}^\varepsilon(t)
 &:=C E_\varepsilon\int_{t-\delta}^{t+\delta}v^\rho(u)\,du.
 \notag
\end{align}
Consequently every error in \eqref{eq:v95-packet-error} has ordered
height at most $C E_\varepsilon\delta$. At each fixed positive
$\rho$ the four zero-residence charges agree, as extended nonnegative
numbers, when $\delta$ decreases after the collision limit.
The common charge is carried entirely by segments having two
nongrazing physical endpoints.
\end{theorem}
\begin{proof}
On a selected segment radial monotonicity gives $z(u)\ge\rho$
throughout. The product rule and the retained full-density estimates give
\[
 |(f_\sigma^\varepsilon)'(u)|\le C E_\varepsilon A_\sigma(u)I_m(u)/z(u).
\]
Absolute continuity, including the one-sided physical traces, implies
\[
 |f_\sigma^\varepsilon(t)-M_\sigma|
 \le\int_{a_\sigma}^{b_\sigma}|(f_\sigma^\varepsilon)'|,
 \qquad |F_\sigma^--F_\sigma^+|
 \le\int_{a_\sigma}^{b_\sigma}|(f_\sigma^\varepsilon)'|.
\]
Moreover $(F_\sigma^--M_\sigma)+(F_\sigma^+-M_\sigma)
=|F_\sigma^--F_\sigma^+|$. All quantities on the latter left side
are nonnegative. Every selected segment lies in
$(t-\delta,t+\delta)$, and its intermediate physical states are
counted once by the original Borel parametrization. Summing the
variation integrals is therefore bounded by the actual $v^\rho$
source on that interval. This proves the finite-count assertions.
Apply \eqref{eq:v94-cutoff} at the previously fixed width $2\delta$.
The common leading constant is independent of the positive $\rho$,
but its finite-count remainder is not being removed at $\rho=0$.

Each packet increases with $\delta$, so its ordered zero-residence
limit exists. The pairwise differences have height tending to zero;
the triangle inequality in both directions proves equality, also for
infinite limits. If either endpoint is grazing, its weighted trace
vanishes, hence $M_\sigma=0$. Restricting the same proof to that
Borel family gives height at most $C E_\varepsilon\delta$ for
its complete original source and for either endpoint packet.
Thus that family's zero-residence charge is zero. No positive lower
residence cutoff was imposed in the finite-count formula.
\end{proof}

\begin{corollary}[A target-gradient version]
\label{cor:v95-target-gradient}
Replace $z_\sigma^-\ge\rho$ in the selected family by
$z_\sigma(t)\ge\rho$. If $0<\delta\le\rho^2/(8\Lambda_\circ)$,
all conclusions of \eqref{eq:v95-packet-error} hold with $v^{\rho/2}$
on the right. In particular the original target-gradient-restricted
short source with a grazing endpoint has zero ordered zero-residence
charge at each fixed positive $\rho$.
\end{corollary}
\begin{proof}
For every earlier point of that segment,
$z(u)^2\ge z(t)^2-2\Lambda_\circ\delta\ge3\rho^2/4$.
At later points $z$ only increases. The complete interval is therefore
above $\rho/2$, and the proof of the theorem applies unchanged.
\end{proof}

The result eliminates all grazing-ended contributions to the
noncritical zero-residence inlet, including segments with a nongrazing
entrance and a grazing exit. It does not estimate the common positive
packet of two nongrazing ends. Nor does equality of the four charges
permit an exchange of the zero-gradient and collision limits.
